\section*{Chapter \ref{ch:mx:child-order-placement} --- Child-Order Placement}

\begin{solution}{exo:mx:child-order-placement:1}
Five lots at half a tick and three at 1.5 ticks: $2.5+4.5=7.0$ ticks for 8 lots, 0.875 ticks a share.
\end{solution}

\begin{solution}{exo:mx:child-order-placement:2}
With $n$ lots ahead the queue loses one at rate $\mu+\theta n$, so the expected time is $\sum_{n=1}^{15}1/(0.69+0.021n)=17.8$ seconds.
\end{solution}

\begin{solution}{exo:mx:child-order-placement:3}
It fills only when the whole bid queue has been consumed or the price has fallen to it (16.7\% of the shares); otherwise the price has usually moved away
by the deadline and the \omterm{def:mx:child-order-placement:risk}{clean-up trade} (83\% of the shares) pays the spread from a worse mid.
\end{solution}

\begin{solution}{exo:mx:child-order-placement:4}
With less time left, a lot that waits is less likely to fill before the deadline, and the \omterm{def:mx:child-order-placement:risk}{clean-up trade} then pays the spread after any adverse move.
The option to wait is worth less, so a smaller predicted risk of an up-move (a lower imbalance) is enough to cross.
\end{solution}

\begin{solution}{exo:mx:child-order-placement:5}
Price moves not caused by a queue emptying (the fundamentalists trade toward a hidden value and move the price whatever the queues), and the
correlation between order flow and those moves: passive buys fill when the price is falling toward them. \omterm{def:mx:the-limit-order-book:orders}{Market orders} larger than a lot, and rates
that depend on the queues, also matter.
\end{solution}

\begin{solution}{exo:mx:child-order-placement:6}
Repricing matters only on the slices where the price ran away: it turns their costly \omterm{def:mx:child-order-placement:risk}{clean-up trades} into later passive fills. Those slices are the
tail of the cost distribution, so removing them shrinks the dispersion (0.72 against 0.89) while the mean moves by little, since most slices fill
either way.
\end{solution}

\begin{solution}{exo:mx:child-order-placement:7}
In 6.0\% of fresh-queue states with the fitted rates, and in 18.6\% with $\lambda$ halved: a thinner, faster-emptying ask makes an up-move likelier,
and crossing at once pays more often.
\end{solution}

\begin{solution}{exo:mx:child-order-placement:8}
The fill is measured against the mid at the fill, not at the decision: the passive order filled because the price came to it, often just before
falling further (adverse selection), and the orders that did not fill were cleaned up after the price moved away. Against the mid at the decision,
joining earned $-0.24$ ticks a share in the simulated market, not $-0.5$.
\end{solution}

\begin{solution}{pb:mx:child-order-placement:1}
\textbf{1.} Choose where and when to post a \omterm{def:mx:the-almgren-chriss-framework:parent}{child order} due by a deadline, and how to revise it, minimising the cost against the mid at the decision.
\textbf{2.} The chance of not being filled by the deadline, and the marketable order that then does the remainder.
\textbf{3.} The spread is one tick 97.7\% of the time: there is no price between the bid and the ask.
\textbf{4.} A one-tick spread 97.7\% of the time, 14.6 lots at the best on average, a mid move every 9.4 seconds, 137 shares a second.
\textbf{5.} Cross 0.70 (0.53), join $-0.24$ (0.89), behind 0.76 (1.59), reprice $-0.22$ (0.72) ticks a share.
\textbf{6.} It walks past the ask into the second level (3.2 lots on average).
\textbf{7.} 91.0\% passively, 9.0\% by the \omterm{def:mx:child-order-placement:risk}{clean-up trade}.
\textbf{8.} It crossed on 6.5\% of the slices and cost 0.03 ticks a share more than repricing.
\textbf{9.} A market sell at rate $\mu$ takes the bid's front lot, lots ahead cancel at $\theta n$, the ask shrinks at $\mu+\theta a$ and grows at
$\lambda$; an empty ask moves the price up a tick and restarts from fresh queues.
\textbf{10.} $V_k=\min(C(r,a),\E[V_{k-1}\mid r,n,a])$, $V_0=C(r,a)$.
\textbf{11.} $\lambda=0.99$, $\theta=0.021$, $\mu=0.69$, a second level of 3.2 lots.
\textbf{12.} It never crosses at the start and expects $-0.40$ ticks a share (crossing: 0.66).
\textbf{13.} 0.86, 0.79 and 0.66.
\textbf{14.} \emph{Named result}: for an 800-share one-minute slice, crossing costs 0.70 ticks a share, joining $-0.24$, joining with repricing
(the plan's choice) $-0.22$, one tick behind 0.76, the imbalance rule $-0.19$; the plan never crosses such a slice at the start, and for a single
lot it crosses above an imbalance of 0.86 with fifteen seconds left, 0.79 with five and 0.66 with three.
\textbf{15.} Waiting's advantage falls from 0.26 ticks below $-0.5$ to nothing measurable above 0.5; crossing is never significantly cheaper.
\textbf{16.} By updating $Q(s,a)$ toward the observed cost plus the best next value on simulated episodes; the state is lots left, lots ahead, lots
at the ask and a five-second time bucket.
\textbf{17.} $-0.39$ ticks a share against $-0.40$, having visited 4.0\% of its table.
\textbf{18.} The real state is larger (more queues, flow, signals), rewards are noisy and episodes costly, so a table cannot be filled.
\textbf{19.} Rest at the best quote and reprice when it moves; cross only small remainders near the deadline when the queues say the price is about
to leave.
\textbf{20.} When little time is left, the order is small, and the queue imbalance says the price is about to move away.
\end{solution}

\begin{solution}{iq:mx:child-order-placement:1}
A \omterm{def:mx:the-limit-order-book:orders}{limit order} at the bid, repriced if the bid moves away and crossed near the deadline: with deep queues and a minute, the fill probability is high
and the half-spread saved is most of the cost.
\iqlookfor{Fill probability versus the spread; a clean-up plan.}
\end{solution}

\begin{solution}{iq:mx:child-order-placement:2}
A deep bid against a thin ask predicts that the ask empties first and the price rises: a resting buy is then likely to be left behind. It lowers the
value of waiting, most when little time is left; crossing is right above a threshold that falls toward the deadline.
\iqlookfor{Prediction of the next move; interaction with time left.}
\end{solution}

\begin{solution}{iq:mx:child-order-placement:3}
State: quantity left, position in the queue, the queues' sizes (and any signal), time left; actions: wait, reprice, cross; cost: fills against the
decision mid plus the terminal clean-up; backward induction on a fitted Markov model of the queues.
\iqlookfor{A complete state; the terminal condition.}
\end{solution}

\begin{solution}{iq:mx:child-order-placement:4}
It loses queue priority each time and multiplies messages (throttles, fees, message-to-trade limits). Add a minimum resting time, a tolerance before
chasing, and state that remembers the queue position it would give up.
\iqlookfor{Priority loss; message costs; hysteresis.}
\end{solution}

\begin{solution}{iq:mx:child-order-placement:5}
Only against a benchmark: the dynamic-programming plan on a fitted model and simple rules, evaluated on the same simulated flow (common random
numbers) and then on live A/B slices. The learner must beat them out of sample, not on the episodes it trained on.
\iqlookfor{Baselines; paired evaluation; out-of-sample.}
\end{solution}

\begin{solution}{iq:mx:child-order-placement:6}
The passive fills are adversely selected (filled when the price was about to go the wrong way), and the 10\% that do not fill are cleaned up after the
price ran away; a colleague crossing when the imbalance warns of a move avoids both.
\iqlookfor{Adverse selection of fills; the cost of the unfilled.}
\end{solution}
